\medskip

Une concession automobile vend deux sortes de véhicules :

\begin{itemize}[label=$\bullet~$]
\item 60\,\% sont des véhicules tout-électrique ;
\item 40\,\% sont des véhicules hybrides rechargeables.
\end{itemize}

\medskip

75\,\% des acheteurs de véhicules tout-électrique et 52\,\% des acheteurs de véhicules hybrides ont la possibilité matérielle d'installer une borne de recharge à domicile.

On choisit un acheteur au hasard et on considère les évènements suivants:

\begin{itemize}[label=$\bullet~$]
\item $E$ : \og l'acheteur choisit un véhicule tout-électrique \fg{} ;
\item $B$ : \og l'acheteur a la possibilité d'installer une borne de recharge à son domicile \fg.
\end{itemize}







On crée un arbre pondéré résumant la situation.

\begin{center}
\bigskip
  \pstree[treemode=R,nodesepA=0pt,nodesepB=4pt,levelsep=3.5cm,nrot=:U]{\TR{}}
 {
 	\pstree[nodesepA=4pt]{\TR{$E$}\naput{$0,6$}}
 	  { 
 		  \TR{$B$}\naput{$0,75$}
 		  \TR{$\overline{B}$}\nbput{\blue $1-0,75=0,25$}	   
 	  }
 	\pstree[nodesepA=4pt]{\TR{$\overline{E}$}\nbput{$0,4$}}
 	  {
 		  \TR{$B$}\naput{$0,52$}
          \TR{$\overline{B}$}\nbput{\blue $1-0,52=0,48$} 
     }
}
\bigskip
\end{center}

\begin{enumerate}
\item La probabilité que l'acheteur choisisse un véhicule tout-électrique et qu'il ait la possibilité d'installer une borne de recharge à son domicile est:\\
$P(E\cap B) = P(E)\times P_E(B)=0,6\times 0,75 = 0,45$.



\item 
D'après la formule des probabilités totales:\\
$P(B)= P(E\cap B) + P \left (\overline{E}\cap B \right ) = 0,45+0,4\times 0,52=0,658$.

\item Un acheteur a la possibilité d'installer une borne de recharge à son domicile.

La probabilité qu'il choisisse un véhicule tout-électrique est:\\
$P_B(E)=\dfrac{P(E\cap B)}{P(B)} = \dfrac{0,45}{0,658}\approx 0,684$.

\item On choisit un échantillon de $20$ acheteurs. On assimile ce prélèvement à un tirage avec remise.
On note $X$ la variable aléatoire qui donne le nombre total d'acheteurs pouvant installer une borne de recharge à leur domicile parmi l'échantillon de $20$~acheteurs.

	\begin{enumerate}
	\item 
On est dans le cas d'une répétition avec remise d'une expérience n'ayant que 2 issues: la possibilité d'installer une borne de recharge, avec la probabilité $p=0,658$, ou non.
Donc la variable aléatoire $X$ qui donne le nombre total d'acheteurs pouvant installer une borne de recharge à leur domicile suit la loi binomiale de paramètres $n=20$ et $p=0,658$.

		\item $P(X = 8) = \ds\binom{20}{8}\times 0,658^8\times (1-0,658)^{20-8} \approx 0,011$

		\item La probabilité qu'au moins 10 acheteurs puissent installer une borne de recharge est:
$P(X\geqslant 10) = 1-P(X\leqslant 9) \approx 1- \np{0,0452} \approx 0,955$

		\item L'espérance de $X$ est: $E(X)=np=20\times 0,658 = 13,16$.

		\item La directrice de la concession décide d'offrir l'installation de la borne de recharge aux acheteurs ayant la possibilité d'en installer une à leur domicile. Cette installation coûte \np{1200}~\euro.
		


L'espérance de $X$ représente le nombre moyen de clients pouvant installer une borne de recharge à leur domicile. L'installation d'une borne coûte \np{1200}~\euro.

Il faut donc prévoir: $13,16\times \np{1200}$ soit \np{15792}~\euro.
	\end{enumerate}
\end{enumerate}

\bigskip
